% Auxiliary proof for the proposed direct integral cycle detector.
% No independent numbered result or main-gate completion is asserted here.
\paragraph{Auxiliary lemma (proved here; awaiting independent audit):
the relative Milnor differential quotient in every characteristic.}
Let $k$ be a field, let $V/k$ be smooth, and let $V\hookrightarrow E$
be a square-zero thickening over $k$ whose ideal $I$ is an invertible
$\mathcal O_V$-module. We use the Zariski sheaf of ordinary Milnor
symbols: on a ring it is the tensor algebra of the group of units,
modulo the two-sided ideal generated by $u\otimes(1-u)$ whenever
both entries are units. Write
\[
 \mathcal J_r=\ker\bigl(\mathcal K_r^M(E)\longrightarrow
                              \mathcal K_r^M(V)\bigr),\qquad r\geq1.
\]
Then this map is surjective, and there is a canonical surjective
homomorphism of sheaves of abelian groups
\[
 \tau_r:\mathcal J_r\longrightarrow
                     I\otimes_{\mathcal O_V}\Omega^{r-1}_{V/k}.
 \tag{RM.1}
\]
Its normalization is
\[
 \tau_r\{1+i,u_2,\ldots,u_r\}
       =i\otimes d\log\bar u_2\wedge\cdots\wedge d\log\bar u_r.
 \tag{RM.2}
\]
It is independent of a local algebra splitting or a local generator
of $I$. If $k$ has characteristic $p>0$, then $p\mathcal J_r=0$,
including when $p=2$. We assert a quotient, not an isomorphism.

\begin{proof}
\emph{Relative generators and exponent.}
For a square-zero surjection of local rings $B\twoheadrightarrow A$
with ideal $I$, the unit map is onto, with kernel $1+I$. The kernel
of the map between their tensor algebras is generated by tensors
having an entry in $1+I$. Moreover each Steinberg relation downstairs
lifts to a Steinberg relation upstairs: if $a,1-a\in A^*$, choose
a lift $u$ of $a$; both $u$ and $1-u$ are units because their
reductions are units. Multiplying these lifted relations by arbitrary
lifted tensors proves that the relative Milnor kernel is generated
by symbols having at least one entry $1+i$, $i\in I$.
This also proves surjectivity. In characteristic $p$ we have
$(1+i)^p=1$, since $I^2=0$ and $pi=0$. Multilinearity therefore kills
each such generator by $p$, in every position. Sheafification
preserves these statements.

\emph{Canonical conormal construction.}
The smoothness of $V$ supplies local $k$-algebra splittings of
$\mathcal O_E\to\mathcal O_V$ by formal smoothness. The lifting
property and the equivalence with smoothness for finitely presented
algebras are \citet[Definition 10.138.1 and Proposition 10.138.13,
Tag 00TH]{frStacksFormalSmoothRelativeMilnor}.
On a splitting, the square-zero algebra is $A\oplus I$, and its
differentials restricted to $A$ are $\Omega^1_{A/k}\oplus I$,
where $i\in I$ corresponds to $di$. Thus the canonical conormal
sequence is short exact:
\[
 0\longrightarrow I\xrightarrow{d}
       \Omega^1_{E/k}\otimes\mathcal O_V
       \longrightarrow\Omega^1_{V/k}\longrightarrow0.
 \tag{RM.3}
\]
This calculation proves injectivity even in characteristic two;
it does not assume that $V\hookrightarrow E$ is a regular immersion.
Taking exterior powers, using that $I$ has rank one, gives canonically
\[
 \ker\left(\Omega^r_{E/k}\otimes\mathcal O_V
                    \longrightarrow\Omega^r_{V/k}\right)
       \simeq I\otimes\Omega^{r-1}_{V/k}.
 \tag{RM.4}
\]
The forward construction sends $i\otimes\omega$ to
$di\wedge\widetilde\omega$. It is independent of the lift
$\widetilde\omega$, because the wedge of two sections of the
rank-one conormal submodule is zero. Its isomorphism property can
be checked on any local splitting of (RM.3).

For any $k$-algebra the map
\[
 \{u_1,\ldots,u_r\}\longmapsto
          d\log u_1\wedge\cdots\wedge d\log u_r
\]
is multilinear and kills the Steinberg relation; indeed
$du\wedge d(1-u)=0$. Restrict this map from $E$ to $V$.
For an element of $\mathcal J_r$ the result lies in the kernel in
(RM.4). This defines (RM.1) without any choice. The identity
$d\log(1+i)|_V=di$ gives (RM.2). Locally every differential $da$
is a multiple of $d\log u$ for a unit $u$: take $u=a$ when $a$
is a unit, and $u=1+a$ otherwise. Products of these expressions
show that (RM.2) generates the target stalks, proving surjectivity.
In degree two a symbol having both entries in $1+I$ maps to zero,
since its restricted differential is a wedge of two conormal
sections. This last fact does not assert that the symbol itself
vanishes in $\mathcal J_2$.
\end{proof}

\paragraph{Compatibility with an actual product of line bundles.}
Let $\mathcal A$ and $\mathcal D$ be line bundles on $E$, and fix an
isomorphism $\mathcal A|_V\simeq\mathcal O_V$. Choose local
trivializations respecting it. The transition functions of
$\mathcal A$ are $1+a_{ij}$, defining $a\in H^1(V,I)$, and let
$g_{ij}$ be those of $\mathcal D$. The symbol cocycle
\[
       \{1+a_{ij},g_{jk}\}
\]
defines a relative class $\xi\in H^2(V,\mathcal J_2)$. Its image
in $H^2(E,\mathcal K_2^M)$ is the ordinary cup product of the two
line-bundle classes. Applying (RM.2) gives exactly
\[
 H^2(\tau_2)(\xi)
       =a\smile d\log(\mathcal D|_V)
            \quad\text{in }H^2(V,I\otimes\Omega^1_{V/k}).
 \tag{RM.5}
\]
This follows on the displayed cocycles and is independent of their
trivializations. In characteristic $p$, if
$H^1(V,\mathcal K_2^M)$ is $p$-divisible, the exact sheaf sequence
\[
 0\longrightarrow\mathcal J_2\longrightarrow\mathcal K_2^M(E)
       \longrightarrow\mathcal K_2^M(V)\longrightarrow0
\]
shows that $H^2(V,\mathcal J_2)\to H^2(E,\mathcal K_2^M)$ is
injective: its kernel is the image of a $p$-divisible group in a
group killed by $p$. A nonzero right side of (RM.5) therefore detects
the nonzero absolute product. This statement uses no comparison
between Milnor and Quillen $K$-theory on the nonreduced $E$.

The same calculation applies to an actual product of $r$ line
bundles, with the first one trivialized on $V$. If the remaining
bundles are $\mathcal D_2,\ldots,\mathcal D_r$, their transition
symbols give $\xi_r\in H^r(V,\mathcal J_r)$ with
\[
 H^r(\tau_r)(\xi_r)=
 a\smile d\log(\mathcal D_2|_V)\smile\cdots\smile
              d\log(\mathcal D_r|_V)
 \quad\text{in }H^r(V,I\otimes\Omega^{r-1}_{V/k}).
 \tag{RM.7}
\]
If $H^{r-1}(V,\mathcal K_r^M)$ is $p$-divisible, then
$H^r(V,\mathcal J_r)\to H^r(E,\mathcal K_r^M)$ is injective
by the identical long-exact-sequence argument. No acyclicity of
the chosen cover for $\mathcal J_r$ is required: the symbol cocycle
maps canonically from \v Cech cohomology to sheaf cohomology, and
its image under $H^r(\tau_r)$ is the displayed coherent cup product.

\paragraph{Independent degree-two Dennis--Stein check.}
There is also a direct degree-two check in relative Quillen
$K$-theory. For $I^2=0$, the Dennis--Stein presentation of
\citet[Chapter III, Definition 5.11, p.~43, and Theorem 5.11.1(b),
p.~44]{frWeibelKbookIII} gives
\[
 \langle i,r\rangle+\langle j,r\rangle
       =\langle i+j,r\rangle \quad(i,j\in I).
\]
Here (D1) changes the entries and (D2) adds them; its correction
term $rij$ is zero. Thus $pK_2(B,I)=0$ when $pI=0$.
The natural homomorphism
\[
 \langle i,r\rangle\longmapsto-i\otimes d\bar r,
 \qquad
 \langle r,i\rangle\longmapsto+i\otimes d\bar r
 \tag{RM.6}
\]
respects (D1)--(D3), including all possible placements of an entry
in $I$: (D2) is additivity and (D3) is the Leibniz identity. For
example, with $s\in I$, (D3) becomes
\[
 st\,d\bar r=-rs\,d\bar t+s\,d(\bar t\bar r).
\]
This is the negative of the Dennis--Stein map in Weibel's
Exercise 5.14(a), pp.~46--47; the sign makes
$\{1+i,u\}\mapsto i\,d\log\bar u$, agreeing with (RM.2).
Nilpotent ideals are contained in the Jacobson radical, as required
by the term ``radical ideal'' in that presentation. This check is
separate from the Milnor proof and is not needed to infer the
nonreduced sheaf comparison.
